\documentclass[12pt,a4paper]{article}

\usepackage[utf8]{inputenc}
\usepackage[T1]{fontenc}
\usepackage[portuguese]{babel}
\usepackage{amsmath,amssymb,amsthm}
\usepackage{geometry}
\usepackage{hyperref}
\usepackage{booktabs}
\usepackage{enumitem}

\geometry{margin=2.4cm}

\hypersetup{
    colorlinks=true,
    linkcolor=blue,
    citecolor=blue,
    urlcolor=blue
}

\newtheorem{teorema}{Teorema}
\newtheorem{proposicao}{Proposição}
\newtheorem{lema}{Lema}
\newtheorem{corolario}{Corolário}
\newtheorem{observacao}{Observação}
\newtheorem{objecao}{Objeção}
\newtheorem{definicao}{Definição}
\newtheorem{problema}{Problema}

\title{
    NUMERIA ES7: \\
    problema de momentos de Hamburger a partir da \\
    fórmula de Guinand--Weil gaussiana (extensão da ES6)
}

\author{}
\date{}

\begin{document}

\maketitle

\begin{abstract}
A ES6 certificou rigorosamente a parcela aritmética da fórmula de
Guinand--Weil para a gaussiana $h_a(t)=e^{-at^2}$ em $a=1/50$, e propôs
uma identidade nova nesta linha, obtida por diferenciação em $a$. Este
documento formaliza a extensão natural desse aparato: a família de
\emph{momentos} $\mu_k(a) = \sum_\rho z_\rho^{2k} e^{-az_\rho^2}$,
suas realizações via as derivadas $k$-ésimas do lado direito da fórmula
explícita, e a matriz de Hankel
$H_n(a) = (\mu_{i+j}(a))_{i,j=0}^{n}$. Mostramos que RH implica
$H_n(a)\succeq 0$ para todo $n,a$, mas que a recíproca \emph{não} vale:
a positividade de Hankel é uma condição \emph{necessária}, não
suficiente. O documento especifica o que é certificável em aritmética
racional exata (valores de $\mu_k$, positividade de $H_n$ em pontos
isolados), o que é equivalente a RH (o limite $n\to\infty$, ou a
positividade para todas as funções-teste de tipo positivo), e o que
permanece aberto. Não há prova de RH. A contribuição é a formalização
de um \emph{teste numérico finito} com conteúdo: uma matriz de Hankel
negativa em algum $(n,a)$ \emph{refutaria} RH, e a trajetória de
$H_n(a)$ codifica informação sobre a distribuição dos zeros.
\end{abstract}

\tableofcontents

\section{Contexto: o que a ES6 estabeleceu}

A ES6, a partir da fórmula global de Guinand--Weil na forma
(Proposition~5 de Carneiro--Chirre--Milinovich, com a convenção de
Fourier $\widehat h(y) = \int_{\mathbb{R}} h(t)e^{-2\pi i t y}\,dt$),
obteve, para $a>0$ e com $z_\rho = (\rho-\tfrac12)/i$ para cada zero
não trivial $\rho$ de $\zeta$, sem pressupor que $z_\rho$ seja real:
\begin{equation}\label{eq:GW}
\boxed{
\sum_{\rho}e^{-az_\rho^2}
=2e^{a/4}+\frac{1}{2\pi}\int_{-\infty}^{\infty}e^{-at^2}\Omega(t)\,dt
-\frac{1}{\sqrt{\pi a}}\sum_{n\ge 2}\frac{\Lambda(n)}{\sqrt n}
e^{-(\log n)^2/(4a)}
}
\end{equation}
onde $\Omega(t) = \operatorname{Re}\psi(\tfrac14+\tfrac{it}{2})-\log\pi$.

A ES6 certificou em aritmética racional exata a parcela positiva
$P(a)=\frac1{\sqrt{\pi a}}\sum_{n\ge2}\frac{\Lambda(n)}{\sqrt n}
e^{-(\log n)^2/(4a)}$ em $a=1/50$:
\[
0{,}0048201252095083525923710170549320614490
\le P(1/50) \le
0.0048201252095083525923710170549320614627,
\]
com cauda $<10^{-40}$. A parcela arquimediana e a soma completa sobre
zeros não foram certificadas com o mesmo rigor.

A ES6 também derivou em $a$, obtendo uma identidade para
$-\sum_\rho z_\rho^2 e^{-az_\rho^2}$. É o ponto de partida desta nota.

\section{A família de momentos}

\begin{definicao}[Momentos gaussianos]
Para $k\ge 0$ e $a>0$, defina
\[
  \mu_k(a) \;=\; \sum_{\rho} z_\rho^{2k}\,e^{-az_\rho^2},
\]
a soma sendo sobre \emph{todos} os zeros não triviais de $\zeta$, com
multiplicidade, e $z_\rho=(\rho-\tfrac12)/i$ possivelmente complexo.
\end{definicao}

\begin{observacao}[Realidade de $\mu_k$]
Se todos os $z_\rho$ são reais (isto é, RH vale), então cada
$z_\rho^{2k} e^{-az_\rho^2} > 0$, e $\mu_k(a)\in\mathbb{R}_{>0}$ para
todo $k,a$. Sob RH, os $\mu_k$ são somas de termos positivos.
\end{observacao}

\begin{proposicao}[$\mu_k$ via diferenciação]\label{prop:deriv}
Suponha que a série em \eqref{eq:GW} converge absolutamente e é
uniformemente diferenciável $k$ vezes em $a$ (o que é consequência do
decaimento gaussiano e da cota clássica $N(T+1)-N(T)\ll\log T$ para a
contagem de zeros). Então, para todo $k\ge 0$ e $a>0$,
\[
  \mu_k(a) \;=\; (-1)^k\,\frac{d^k}{da^k}\,\mathrm{RHS}_0(a),
\]
onde $\mathrm{RHS}_0(a)$ denota o lado direito de \eqref{eq:GW}.
\end{proposicao}

\begin{proof}
Para $k=0$ é a identidade \eqref{eq:GW}. Para $k\ge 1$, diferencia-se
$k$ vezes em ambos os lados; a troca de soma e derivada é válida
porque $|z_\rho|^{2k}e^{-a z_\rho^2}$ decai gaussianamente e a densidade
de zeros é $O(\log T)$. A fórmula explícita para $\mu_k$ em termos de
$\Lambda$, $\Omega$ e suas derivadas é obtida pela diferenciação
termo a termo do lado direito de \eqref{eq:GW}.
\end{proof}

\begin{observacao}[Certificabilidade em $\mathbb{Q}$]
Para $k$ fixo e $a$ fixo, o valor de $\mu_k(a)$ no lado direito de
\eqref{eq:GW} é uma soma de termos: potências de primos (finitas até
$N$, mais cauda majorada) e integrais de $\Omega^{(j)}$ contra $t^{2j}
e^{-at^2}$. Tudo isso é certificável pelo aparato racional da ES6. A
única parte que a ES6 não certifica é a quadratura da digama
$\psi$; mas isso é uma limitação \emph{técnica}, não \emph{estrutural}.
\end{observacao}

\section{A matriz de Hankel}

\begin{definicao}[Matriz de Hankel de momentos]\label{def:hankel}
Para $n\ge 0$ e $a>0$, defina
\[
  H_n(a) \;=\; \bigl(\mu_{i+j}(a)\bigr)_{i,j=0}^{n}
  \in\; \mathbb{R}^{(n+1)\times(n+1)}.
\]
\end{definicao}

\begin{teorema}[RH implica positividade de Hankel]\label{thm:RH-PSD}
Se a Hipótese de Riemann é verdadeira, então para todo $n\ge 0$ e todo
$a>0$, $H_n(a)$ é semidefinida positiva, $H_n(a)\succeq 0$.
\end{teorema}

\begin{proof}
Sob RH, todos os $z_\rho$ são reais. Seja $\nu_a$ a medida positiva em
$\mathbb{R}$ dada por
\[
  \nu_a \;=\; \sum_{\rho}\delta_{z_\rho^2}\cdot e^{-az_\rho^2}
\]
(medida atômica com massa $e^{-az_\rho^2}$ no ponto $z_\rho^2\ge 0$).
Então
\[
  \mu_k(a) \;=\; \sum_\rho (z_\rho^2)^k\,e^{-az_\rho^2}
  \;=\; \int_{\mathbb{R}} x^k\,d\nu_a(x),
\]
isto é, $\mu_k(a)$ é o $k$-ésimo momento da medida positiva $\nu_a$
em $\mathbb{R}$. A matriz de Hankel $H_n(a)$ é a matriz de Gram dos
monômios $\{1,x,x^2,\dots,x^n\}$ em $L^2(\mathbb{R},\nu_a)$, logo
$H_n(a) = M_n^*M_n \succeq 0$, onde $M_n$ tem colunas $x^j$ em
$L^2(\nu_a)$.
\end{proof}

\begin{corolario}[Teste de refutação]\label{cor:refute}
Se existirem $n\ge 0$ e $a>0$ tais que $H_n(a)\not\succeq 0$, então RH
é falsa.
\end{corolario}

\begin{proof}
Contrapositiva do Teorema \ref{thm:RH-PSD}.
\end{proof}

\section{O que a recíproca daria — e por que não dá}

A recíproca do Teorema \ref{thm:RH-PSD} seria a afirmação:
\begin{center}
\fbox{
\begin{minipage}{0.85\textwidth}
\emph{(Recíproca falsa desejada.)} Se $H_n(a)\succeq 0$ para todo
$n\ge 0$ e todo $a>0$, então RH é verdadeira.
\end{minipage}
}
\end{center}

Esta recíproca é \emph{falsa} como enunciada, por três razões
independentes e cumulativas.

\begin{objecao}[Uma única gaussiana não é suficiente]
A positividade de $H_n(a)$ para um $a$ fixo diz que a sequência
$(\mu_k(a))_k$ é a sequência de momentos de \emph{alguma} medida
positiva em $\mathbb{R}$ (por Hamburger), mas não necessariamente da
medida $\nu_a = \sum_\rho \delta_{z_\rho^2} e^{-az_\rho^2}$
\emph{original}. A positividade da forma de Weil exige a família
completa de funções-teste de tipo positivo, não apenas as gaussianas.
\end{objecao}

\begin{objecao}[O limite $n\to\infty$ é intratável]
O teorema de Hamburger caracteriza as sequências de momentos de medidas
positivas em $\mathbb{R}$ via \emph{todas} as matrizes de Hankel
$H_n\succeq 0$ para todo $n$. Verificar isso para \emph{todo} $n$ é,
na prática, um problema de controle uniforme de $\mu_k(a)$ em $k$ e
$a$, o que envolve as mesmas dificuldades de estimativa da função de
Mertens que caracterizam o obstáculo original do M1.
\end{objecao}

\begin{objecao}[A medida $\nu_a$ não é única]
Mesmo que todas as $H_n(a)$ sejam PSD para todo $n$ e todo $a$, isso
caracteriza uma \emph{classe} de medidas positivas com os mesmos
momentos. Se a medida $\nu_a$ original (que depende dos zeros de
$\zeta$) pertence ou não a essa classe não é determinado apenas pelos
momentos; é um problema separado.
\end{objecao}

\begin{observacao}[Formulação correta do que seria prova]
A recíproca verdadeira, sob a qual a positividade de Hankel \emph{seria}
suficiente, exigiria identificar a medida limite
$\nu = \lim_{a\to 0}\nu_a$ com a medida de contagem dos $z_\rho^2$ e
verificar positividade contra \emph{todas} as funções-teste de tipo
positivo de Weil, não apenas gaussianas. Isso é exatamente a forma de
Weil, que é equivalente a RH; portanto a recíproca, na sua forma útil,
\emph{não simplifica} o problema.
\end{observacao}

\section{O que é certificável — e por que tem conteúdo}

Apesar da recíproca falhar, o Teorema \ref{thm:RH-PSD} e o Corolário
\ref{cor:refute} têm conteúdo computacional genuíno.

\begin{proposicao}[Teste de PSD em $\mathbb{Q}$]\label{prop:PSD-Q}
Fixem $n\ge 0$, $a\in\mathbb{Q}_{>0}$ e um inteiro $N\ge 64$. Então,
usando o aparato racional da ES6 para certificar $\mu_0,\dots,\mu_{2n}$
até $N$ com cauda majorada, é possível decidir com certificado
racional se a matriz truncada $\widetilde H_n^{(N)}(a)$ satisfaz
$\widetilde H_n^{(N)}(a)\succeq \lambda I$ para algum
$\lambda\in\mathbb{Q}_{>0}$ explicitamente calculável.
\end{proposicao}

\begin{proof}[Esboço]
Cada $\mu_k(a)$ é uma soma finita de termos certificáveis em intervalos
racionais (potências de primos até $N$, mais cauda majorada) mais uma
integral de $\Omega^{(j)}$ contra um peso gaussiano. A certificação
intervalar da digama é a única lacuna técnica; supondo-a resolvida
(por exemplo, via biblioteca Arb/FLINT), os intervalos racionais dos
$\mu_k$ propagam-se para $\widetilde H_n^{(N)}(a)$ via aritmética de
intervalos, e um teste de Cholesky racional certifica positividade ou
encontra um vetor $v\in\mathbb{Q}^{n+1}$ tal que
$v^\top\widetilde H_n^{(N)}(a)v < 0$.
\end{proof}

\begin{observacao}[Três desfechos possíveis do teste]
\begin{enumerate}[label=(\roman*)]
\item \textbf{PSD confirmado} (com $N$ grande o suficiente): evidência
numérica de que $\mu_0,\dots,\mu_{2n}$ são compatíveis com RH. Resultado
modesto mas honesto.
\item \textbf{Não-PSD confirmado}: \emph{RH é falsa}. Como não
esperamos isso, esse desfecho é uma falsificação improvável, mas o teste
é o que daria o certificado.
\item \textbf{Inconclusivo} (intervalos dos $\mu_k$ muito largos para
decidir): sinal de que $N$ ou a precisão precisam aumentar.
\end{enumerate}
\end{observacao}

\section{Extensão concreta do aparato ES6}

\begin{problema}[ES7-A: certificar $\mu_2, \mu_3, \mu_4$]\label{prob:mu234}
Estender a certificação racional da ES6 a $\mu_k(a)$ para
$k=2,3,4$, $a\in\{1/50, 1/25, 1/10\}$, $N\ge 128$, com cauda
majorada analiticamente. Aderir à estrutura de arredondamento para
fora em escala $10^{40}$ já usada.
\end{problema}

\begin{problema}[ES7-B: montar e testar $H_3(a)$]\label{prob:H3}
Para $n\in\{1,2,3\}$ e $a\in\{1/50, 1/25, 1/10\}$, computar
$\widetilde H_n^{(N)}(a)$ com intervalos racionais para $\mu_0,\dots,\mu_{2n}$,
e rodar Cholesky racional para certificar $\widetilde H_n^{(N)}(a)\succeq \lambda I$
com $\lambda\in\mathbb{Q}_{>0}$ explícito, ou detectar não-positividade.
\end{problema}

\begin{problema}[ES7-C: trajetória de $H_n(a)$]\label{prob:traj}
Calcular o menor autovalor $\lambda_{\min}(H_n(a))$ para
$n\in\{1,2,3,4\}$ e $a\in\{1/5, 1/10, 1/20, 1/50, 1/100, 1/500\}$,
plotar $\lambda_{\min}$ contra $\log a$ e contra $n$, e comparar
com a previsão sob RH (todos positivos) e sob uma distribuição
GUE de zeros na faixa crítica (eventuais negativos). Se as duas
previsões divergem em escala, a estrutura de Hankel detecta algo.
\end{problema}

\begin{problema}[ES7-D: certificação intervalar da digama]\label{prob:digamma}
Implementar quadratura intervalar rigorosa para
$\int_{-\infty}^{\infty} t^{2j} e^{-at^2}\Omega^{(j)}(t)\,dt$
usando Arb/FLINT ou equivalente, com limitação de cauda e
arredondamento explícito. Isso remove a única lacuna técnica apontada
na Proposição \ref{prop:PSD-Q}.
\end{problema}

\section{O que a linha ES6/ES7 faz e não faz}

\begin{center}
\renewcommand{\arraystretch}{1.35}
\begin{tabular}{p{6.5cm} p{8.5cm}}
\toprule
\textbf{Afirmação} & \textbf{Status} \\
\midrule
Fórmula de Guinand--Weil para gaussiana $h_a$ & Clássica, incondicional. \\
Identidade derivada em $a$ (ES6) & Correta, não nova na literatura,
nova nesta linha. \\
Certificação de $P(1/50)$ com cauda $<10^{-40}$ & Correta, verificada
em $\mathbb{Q}$. \\
Extensão a $\mu_k(a)$ para $k\ge 2$ & Direta, viável (Problema
\ref{prob:mu234}). \\
$H_n(a)\succeq 0$ sob RH & Demonstrado (Teorema \ref{thm:RH-PSD}). \\
$H_n(a)\not\succeq 0$ $\Rightarrow$ RH falsa & Demonstrado (Corolário
\ref{cor:refute}). \\
Recíproca: $H_n\succeq 0\;\forall n,a\Rightarrow$ RH & \textbf{Falsa}
(Objeções 1--3). \\
Prova de RH & \textbf{Não obtida e não visível neste caminho}. \\
\bottomrule
\end{tabular}
\end{center}

\section{Onde a matemática verdadeiramente diferente está}

A linha de momentos de Hankel é \emph{honesta}: produz certificados
finitos com conteúdo (Corolário \ref{cor:refute}) e um programa
experimental coerente (Problemas \ref{prob:mu234}--\ref{prob:traj}).
Mas ela \emph{não é} a matemática diferente que fecharia RH. As
tentativas sérias na literatura (não desta linha) seguem três rotas
estruturalmente distintas:

\begin{enumerate}[label=(\roman*)]
\item \textbf{Geometria aritmética sobre $\mathbb{F}_q$.} A RH é
\emph{teorema} para curvas sobre corpos finitos (Weil, Deligne). A
prova usa cohomologia étale e positividade geométrica do produto de
interseção. Transportar essa positividade para
$\operatorname{Spec}\mathbb{Z}$ é o programa de Connes (cohomologia
Weil--étale), Scholze (prismas), e outros. Nenhuma prova existe;
é o único caminho onde positividade é provada em análogo estrutural.
\item \textbf{Espectro de um operador autoadjunto explícito.}
Hilbert--Pólya continua sendo a única reformulação em que
``positividade'' vira ``autoadjunção''. A ES6 já mostrou (via
experimento adversarial ES4) que acertar a contagem assintótica
\emph{não basta}; é preciso um operador cujo \emph{espectro exato}
seja $\{z_\rho\}$. Nenhum operador explícito conhecido tem essa
propriedade.
\item \textbf{Formulação adélica de Bombieri.} A fórmula explícita de
Weil tem versão adélica onde o lado arquimediano é substituído por
um produto sobre $p$-ádicos. A positividade é mais rica no contexto
adélico e já produziu teoremas parciais. A ES6/ES7 está na versão
\emph{real}, a mais fraca.
\end{enumerate}

\begin{observacao}
Este documento não apresenta uma prova da Hipótese de Riemann. Ele
formaliza o problema de momentos de Hamburger como extensão natural
da ES6, demonstra o que é demonstrável (Teorema \ref{thm:RH-PSD} e
Corolário \ref{cor:refute}), refuta a recíproca com três objeções
independentes, e especifica um programa computacional concreto
(Problemas \ref{prob:mu234}--\ref{prob:digamma}) com o mesmo rigor
da ES6. Nenhuma dessas etapas aproxima uma prova; todas aproximam
uma \emph{compreensão} de onde a prova não está.
\end{observacao}

\end{document}